\begin{definition}[Parent commuting Hamiltonian on the tripartite space]
    \label{def:tripartite_parent_commuting_hamiltonian}
    \lean{TripartiteDecorrelation.HasGroundSpaceIntersection,
      TripartiteDecorrelation.CommutingParentHamiltonian,
      TripartiteDecorrelation.HasCommutingParentHamiltonian}
    \leanok
    \uses{def:tripartite_local_lifts}
    There are orthogonal projectors $P_{AX}$ and $P_{XB}$ on the two local
    factors whose lifts commute and whose ranges satisfy the source
    ground-space intersection condition
    \[
        \ran P
        =\ran\widehat P_{AX}
          \cap\ran\widehat P_{XB}.
    \]
    Thus $Q_{AX}=\mathbf 1-P_{AX}$ and
    $Q_{XB}=\mathbf 1-P_{XB}$ are exactly the commuting parent terms of
    \cite[Definition~D.2, lines 2205--2218]{Cirac2016MPDO_arXiv}.
\end{definition}

\begin{theorem}[Ground-space intersection and projector product]
    \label{thm:tripartite_ground_space_intersection_iff_product}
    \lean{TripartiteDecorrelation.range_mul_eq_inter,
      TripartiteDecorrelation.hermitian_idempotent_eq_of_range_eq,
      TripartiteDecorrelation.hasGroundSpaceIntersection_iff_product_eq,
      TripartiteDecorrelation.CommutingParentHamiltonian.hproduct}
    \leanok
    \uses{def:tripartite_parent_commuting_hamiltonian}
    Let $P$, $P_{AX}$, and $P_{XB}$ be orthogonal projectors, and write their
    lifts to the tripartite space as
    \[
        \widehat P_{AX}=P_{AX}\otimes\mathbf 1_B,
        \qquad
        \widehat P_{XB}=\mathbf 1_A\otimes P_{XB}.
    \]
    Suppose that the lifted projectors commute:
    \[
        \widehat P_{AX}\widehat P_{XB}
        =\widehat P_{XB}\widehat P_{AX}.
    \]
    Then
    \[
        \ran P
        =\ran\widehat P_{AX}\cap
          \ran\widehat P_{XB}
        \quad\Longleftrightarrow\quad
        \widehat P_{AX}\widehat P_{XB}=P.
    \]
    This is the passage from
    \cite[Definition~D.2, lines 2205--2218]{Cirac2016MPDO_arXiv} to the
    projector identity at lines 2279--2289.
\end{theorem}

\begin{proof}
    \leanok
    A vector lies in the range of
    $\widehat P_{AX}\widehat P_{XB}$ precisely when it is fixed by both
    commuting idempotents. Hence
    \[
        \ran(\widehat P_{AX}\widehat P_{XB})
        =\ran\widehat P_{AX}\cap
          \ran\widehat P_{XB}.
    \]
    The product is again an orthogonal projector.  Orthogonal projectors with
    equal ranges are equal, giving the forward implication; the reverse
    implication follows by substituting the product identity.
\end{proof}

\begin{theorem}[Parent commuting Hamiltonian--complex-linear decorrelation equivalence]
    \label{thm:tripartite_parent_iff_complex_linear_decorrelated}
    \lean{TripartiteDecorrelation.CommutingParentHamiltonian.isDecorrelated,
      TripartiteDecorrelation.IsDecorrelated.hasCommutingParentHamiltonian,
      TripartiteDecorrelation.parentHamiltonian_iff_decorrelated}
    \leanok
    \uses{def:tripartite_decorrelation,
      def:tripartite_parent_commuting_hamiltonian,
      thm:tripartite_support_products,
      thm:tripartite_ground_space_intersection_iff_product}
    Let $P=P^\dagger=P^2$.  Then $P$ is the ground-space projector of a
    parent commuting Hamiltonian on $AX$ and $XB$ if and only if regions
    $A$ and $B$ satisfy complex-linear decorrelation.
\end{theorem}

\begin{proof}
    \leanok
    Starting from decorrelation, take the support projectors of
    $\rho_{AX}$ and $\rho_{XB}$.  The preceding product theorem proves their
    commutation and identifies their product with $P$.

    Starting from commuting parent projectors, first use the ground-space
    intersection theorem to write
    $P=\widehat P_{AX}\widehat P_{XB}
      =\widehat P_{XB}\widehat P_{AX}$.  Locality gives
    $[O_A,\widehat P_{XB}]=[O_B,\widehat P_{AX}]=0$, while
    \[
        \widehat P_{XB}(\mathbf 1-P)\widehat P_{AX}=0.
    \]
    Therefore, using the two locality commutators,
    \[
    \begin{aligned}
        P O_A(\mathbf 1-P)O_B P
        &=(\widehat P_{AX}\widehat P_{XB})O_A(\mathbf 1-P)O_B
          (\widehat P_{AX}\widehat P_{XB})\\
        &=\widehat P_{AX}O_A
          \bigl[\widehat P_{XB}(\mathbf 1-P)\widehat P_{AX}\bigr]
          O_B\widehat P_{XB}\\
        &=0.
    \end{aligned}
    \]
\end{proof}

\begin{theorem}[Parent commuting Hamiltonian--observable decorrelation equivalence]
    \label{thm:tripartite_parent_iff_decorrelated}
    \lean{TripartiteDecorrelation.parentHamiltonian_iff_observableDecorrelated}
    \leanok
    \uses{def:tripartite_observable_decorrelation,
      thm:tripartite_observable_decorrelation_iff,
      thm:tripartite_parent_iff_complex_linear_decorrelated}
    Let $P=P^\dagger=P^2$.  Then $P$ is the ground-space projector of a
    parent commuting Hamiltonian on $AX$ and $XB$ if and only if regions
    $A$ and $B$ are decorrelated when tested on Hermitian observables.
    This is \cite[Proposition~D.3, lines 2221--2289]{Cirac2016MPDO_arXiv}.
\end{theorem}

\begin{proof}
    \leanok
    Combine the complex-linear parent equivalence with the decomposition
    $O=H+\mathrm{i}K$ from
    Theorem~\ref{thm:tripartite_observable_decorrelation_iff}.
\end{proof}

\begin{definition}[Decorrelation]\label{def:is_decorrelated}
    \lean{Decorrelation.IsDecorrelated}
    \leanok
    Given an endomorphism $P_K$ and sets of observables $\mathcal{O}_A$,
    $\mathcal{O}_B$, regions $A$ and $B$ are \emph{decorrelated} w.r.t.\ $K$
    when
    $P_K \circ O_A \circ (1 - P_K) \circ O_B \circ P_K = 0$
    for all $O_A \in \mathcal{O}_A$, $O_B \in \mathcal{O}_B$.
\end{definition}

\begin{definition}[Idempotent-product parent-commuting structure]\label{def:has_commuting_parent_ham}
    \lean{Decorrelation.HasCommutingParentHam}
    \leanok
    The structure used below records only the idempotent-product
    consequence of the parent commuting Hamiltonian condition of
    \cite[Appendix~D.2]{Cirac2016MPDO_arXiv}. It consists of commuting
    idempotent endomorphisms $P_{AX}$ and $P_{XB}$ such that, for the
    idempotent $P_K$ associated to $K$,
    \[
        P_{AX}\circ P_{XB}=P_K.
    \]
    In the source, this is the equation obtained from local projectors
    $Q_{AX}$ and $Q_{XB}$ satisfying $[Q_{AX},Q_{XB}]=0$ and
    $K_{AXB}=(K_{AX}\otimes H_B)\cap(H_A\otimes K_{XB})$.
    % Scope restriction: this definition records only the idempotent-product
    % part of the source condition. See
    % docs/paper-gaps/cpsv16_parent_commuting_hamiltonian_scope.tex.
\end{definition}

\begin{theorem}[Idempotence gives tautological idempotent-product data]%
    \label{thm:has_commuting_parent_ham_of_idem}
    \lean{Decorrelation.HasCommutingParentHam.ofIdem}
    \leanok
    \uses{def:has_commuting_parent_ham}
    If
    \[
        P_K \circ P_K = P_K,
    \]
    then $P_K$ has idempotent-product data.
\end{theorem}

\begin{proof}
    \leanok
    \uses{def:has_commuting_parent_ham}
    Take $P_{AX} = P_{XB} = P_K$. Their commutativity is immediate, and the
    product identity is exactly the assumed idempotence of $P_K$.
\end{proof}

\begin{lemma}[Cancellation for commuting idempotents]
    \label{lem:comp_complement_comm_zero}
    \lean{LinearMap.comp_complement_comm_zero}
    \leanok
    If $P$ and $Q$ are commuting idempotents, then
    \[
        Q \circ (1 - P \circ Q) \circ P = 0.
    \]
\end{lemma}

\begin{proof}
    \leanok
    The cancellation is the computation
    \[
        Q(1-PQ)P
        =
        QP-QPQP
        =
        QP-QPPQ
        =
        QP-QP
        =
        0,
    \]
    using commutativity and idempotence.
\end{proof}

\begin{theorem}[Idempotence of the product projector]%
    \label{thm:pk_idem}
    \lean{Decorrelation.HasCommutingParentHam.pK_idem}
    \leanok
    \uses{def:has_commuting_parent_ham}
    If $P_K = P_{AX} \circ P_{XB}$ is idempotent-product data, then
    \[
        P_K \circ P_K = P_K.
    \]
\end{theorem}

\begin{proof}
    \leanok
    \uses{def:has_commuting_parent_ham}
    Since $P_{AX}$ and $P_{XB}$ are idempotent and commute, their product
    $P_{AX} \circ P_{XB}$ is idempotent. Substituting
    $P_K = P_{AX} \circ P_{XB}$ gives the claim.
\end{proof}

\begin{theorem}[Absorption by the left projector]%
    \label{thm:pAX_comp_pK}
    \lean{Decorrelation.HasCommutingParentHam.pAX_comp_pK}
    \leanok
    \uses{def:has_commuting_parent_ham}
    If $P_K = P_{AX} \circ P_{XB}$ is idempotent-product data, then
    \[
        P_{AX} \circ P_K = P_K.
    \]
\end{theorem}

\begin{proof}
    \leanok
    \uses{def:has_commuting_parent_ham}
    Using $P_K = P_{AX} \circ P_{XB}$ and the idempotence of $P_{AX}$,
    \[
        P_{AX} \circ P_K
        = P_{AX} \circ P_{AX} \circ P_{XB}
        = P_{AX} \circ P_{XB}
        = P_K.
    \]
\end{proof}

\begin{theorem}[Absorption by the right projector]%
    \label{thm:pK_comp_pXB}
    \lean{Decorrelation.HasCommutingParentHam.pK_comp_pXB}
    \leanok
    \uses{def:has_commuting_parent_ham}
    If $P_K = P_{AX} \circ P_{XB}$ is idempotent-product data, then
    \[
        P_K \circ P_{XB} = P_K.
    \]
\end{theorem}

\begin{proof}
    \leanok
    \uses{def:has_commuting_parent_ham}
    Using $P_K = P_{AX} \circ P_{XB}$ and the idempotence of $P_{XB}$,
    \[
        P_K \circ P_{XB}
        = P_{AX} \circ P_{XB} \circ P_{XB}
        = P_{AX} \circ P_{XB}
        = P_K.
    \]
\end{proof}

\begin{theorem}[Cross-absorption by the right projector]%
    \label{thm:pXB_comp_pK}
    \lean{Decorrelation.HasCommutingParentHam.pXB_comp_pK}
    \leanok
    \uses{def:has_commuting_parent_ham}
    If $P_K = P_{AX} \circ P_{XB}$ is idempotent-product data, then
    \[
        P_{XB} \circ P_K = P_K.
    \]
\end{theorem}

\begin{proof}
    \leanok
    \uses{def:has_commuting_parent_ham}
    By commutativity,
    $P_{XB} \circ P_K = P_{XB} \circ P_{AX} \circ P_{XB}
        = P_{AX} \circ P_{XB} \circ P_{XB}$, and idempotence of $P_{XB}$
    yields $P_{XB} \circ P_K = P_K$.
\end{proof}

\begin{theorem}[Cross-absorption by the left projector]%
    \label{thm:pK_comp_pAX}
    \lean{Decorrelation.HasCommutingParentHam.pK_comp_pAX}
    \leanok
    \uses{def:has_commuting_parent_ham}
    If $P_K = P_{AX} \circ P_{XB}$ is idempotent-product data, then
    \[
        P_K \circ P_{AX} = P_K.
    \]
\end{theorem}

\begin{proof}
    \leanok
    \uses{def:has_commuting_parent_ham}
    By commutativity,
    $P_K \circ P_{AX} = P_{AX} \circ P_{XB} \circ P_{AX}
        = P_{AX} \circ P_{AX} \circ P_{XB}$, and idempotence of $P_{AX}$
    yields $P_K \circ P_{AX} = P_K$.
\end{proof}

\begin{theorem}[Reverse product identity]%
    \label{thm:reverse_product}
    \lean{Decorrelation.HasCommutingParentHam.reverse_product}
    \leanok
    \uses{def:has_commuting_parent_ham}
    If $P_K = P_{AX} \circ P_{XB}$ is idempotent-product data, then
    \[
        P_{XB} \circ P_{AX} = P_K.
    \]
\end{theorem}

\begin{proof}
    \leanok
    \uses{def:has_commuting_parent_ham}
    This follows from commutativity of $P_{AX}$ and $P_{XB}$ together
    with the identity $P_{AX} \circ P_{XB} = P_K$.
\end{proof}

\begin{theorem}[Commuting complementary projectors]%
    \label{thm:complement_comm}
    \lean{Decorrelation.HasCommutingParentHam.complement_comm}
    \leanok
    \uses{def:has_commuting_parent_ham}
    If $P_K = P_{AX} \circ P_{XB}$ is idempotent-product data and
    \[
        Q_{AX} = 1 - P_{AX}, \qquad Q_{XB} = 1 - P_{XB},
    \]
    then
    \[
        Q_{AX} \circ Q_{XB} = Q_{XB} \circ Q_{AX}.
    \]
\end{theorem}

\begin{proof}
    \leanok
    \uses{def:has_commuting_parent_ham}
    Expanding both products and using $P_{AX} \circ P_{XB}
        = P_{XB} \circ P_{AX}$ shows that the complementary projectors commute.
\end{proof}

\begin{theorem}[Idempotent-product data imply decorrelation]%
    \label{thm:commuting_ham_is_decorrelated}
    \lean{Decorrelation.commutingHam_isDecorrelated}
    \leanok
    \uses{def:is_decorrelated, def:has_commuting_parent_ham}
    If $P_{AX}$ and $P_{XB}$ are idempotent endomorphisms satisfying
    $P_{AX} \circ P_{XB} = P_K$ and
    $P_{AX} \circ P_{XB} = P_{XB} \circ P_{AX}$, and
    $A$-observables commute with $P_{XB}$ while $B$-observables commute
    with $P_{AX}$, then regions $A$ and $B$ are decorrelated w.r.t.\ $K$.
    This is the idempotent-product form of the backward direction of
    \cite[Appendix~D, Section~D.2, Proposition~D.3]{Cirac2016MPDO_arXiv}.
\end{theorem}

\begin{proof}
    \leanok
    \uses{lem:comp_complement_comm_zero}
    The key identity is $P_{XB} \circ (1 - P_{AX} \circ P_{XB}) \circ P_{AX} = 0$.
    Substituting $P_K = P_{AX} \circ P_{XB}$ and using
    $O_A \circ P_{XB} = P_{XB} \circ O_A$ and
    $O_B \circ P_{AX} = P_{AX} \circ O_B$ gives
    \[
        P_K \circ O_A \circ (1 - P_K) \circ O_B \circ P_K
        = P_{AX} \circ O_A \circ
          \bigl[P_{XB} \circ (1 - P_{AX} \circ P_{XB}) \circ P_{AX}\bigr]
          \circ O_B \circ P_{XB} = 0.
    \]
\end{proof}

\begin{theorem}[Decorrelating idempotent-product data]%
    \label{thm:has_commuting_parent_ham_is_decorrelated}
    \lean{Decorrelation.HasCommutingParentHam.isDecorrelated}
    \leanok
    \uses{thm:commuting_ham_is_decorrelated, def:has_commuting_parent_ham, def:is_decorrelated}
    Corollary of Theorem~\ref{thm:commuting_ham_is_decorrelated}: if
    $P_K$ has idempotent-product data, the $A$-observables
    commute with $P_{XB}$, and the $B$-observables commute with $P_{AX}$, then
    regions $A$ and $B$ are decorrelated w.r.t.\ $K$.
\end{theorem}

\begin{proof}
    \leanok
    \uses{thm:commuting_ham_is_decorrelated}
    Apply Theorem~\ref{thm:commuting_ham_is_decorrelated} to the witnessing
    endomorphisms $P_{AX}$ and $P_{XB}$ supplied by the idempotent-product data.
\end{proof}

\begin{lemma}[Frustration-free identity]\label{lem:frustration_free_ham_eq}
    \lean{LinearMap.frustration_free_ham_eq}
    \leanok
    For any endomorphisms $P$ and $Q$,
    \[
        (1-P) + (1-Q) - (1-P)\circ(1-Q) = 1 - P\circ Q.
    \]
\end{lemma}

\begin{proof}
    \leanok
    Expanding the left-hand side and cancelling gives the result.
\end{proof}

\begin{theorem}[Idempotent-product Hamiltonian identity]%
    \label{thm:hamiltonian_eq}
    \lean{Decorrelation.HasCommutingParentHam.hamiltonian_eq}
    \leanok
    \uses{def:has_commuting_parent_ham, lem:frustration_free_ham_eq}
    If $P_K = P_{AX} \circ P_{XB}$ is idempotent-product data, then
    \[
        Q_{AX} + Q_{XB} - Q_{AX} \circ Q_{XB} = 1 - P_K,
    \]
    where $Q_{AX} = 1 - P_{AX}$ and $Q_{XB} = 1 - P_{XB}$ are the
    complementary projectors.
    This is the algebraic identity obtained from the complementary projectors
    in \cite[Appendix~D, Section~D.2]{Cirac2016MPDO_arXiv}; it is not the full
    source parent commuting Hamiltonian condition.
\end{theorem}

\begin{proof}
    \leanok
    \uses{lem:frustration_free_ham_eq}
    Applying Lemma~\ref{lem:frustration_free_ham_eq} to $P_{AX}$ and
    $P_{XB}$, and substituting $P_{AX} \circ P_{XB} = P_K$.
\end{proof}

\begin{theorem}[Ground-space membership]\label{thm:mem_ground_iff}
    \lean{Decorrelation.HasCommutingParentHam.mem_ground_iff}
    \leanok
    \uses{def:has_commuting_parent_ham}
    If $P_K = P_{AX} \circ P_{XB}$ is idempotent-product data, then
    \[
        P_K v = v
        \quad\Longleftrightarrow\quad
        P_{AX} v = v \text{ and }  P_{XB} v = v.
    \]
    This is the algebraic consequence corresponding to equation~(D.2) from
    \cite{Cirac2016MPDO_arXiv}:
    $K_{AXB} = (K_{AX} \otimes H_B) \cap (H_A \otimes K_{XB})$.
    The source equation itself also carries the tensor-product local-subspace
    hypotheses.
\end{theorem}

\begin{proof}
    \leanok
    \uses{def:has_commuting_parent_ham, thm:pAX_comp_pK, thm:pXB_comp_pK}
    ($\Rightarrow$) If $P_K v = v$, then $P_{AX}(P_K v) = (P_{AX} \circ P_K) v
        = P_K v = v$ by left absorption, and similarly for $P_{XB}$.

    ($\Leftarrow$) If $P_{AX} v = v$ and $P_{XB} v = v$, then
    $P_K v = (P_{AX} \circ P_{XB}) v = P_{AX}(P_{XB} v) = P_{AX} v = v$.
\end{proof}
